%% =====================================================================
%%  B6-T-03-03 -- what B6 contributes to "The Covert Aggressor"
%%  -------------------------------------------------------------------
%%  LAYER A: APPEND-ONLY. Written once at the entry's write, then left
%%  alone. If the source has more to say later, add a bullet -- do not
%%  rewrite. Material found ONLY here goes in the `unique` slot with
%%  @unique set to yes, which makes it undeletable (rule R10).
%% =====================================================================
%% @kind: source
%% @id: B6-T-03-03
%% @book: B6
%% @topic: T-03-03
%% @locator: chs. 1--6, esp. pp. 9--15, 30--34, 62--66
%% @status: distilled
%% @unique: yes
%% @integrated: yes
%% @updated: 2026-09-19

\dossier{B6}{T-03-03}{\bookshort{B6}}{chs. 1--6, esp. pp. 9--15, 30--34, 62--66}

\begin{distilled}
  \item Covert-aggressive personalities fight hard for everything they want but conceal their aggressive intentions; charming and genial on the surface, calculating and ruthless underneath.
  \item B6's clinical corrective: the traditional lens reads everyone as neurotic --- hung up on fears and defences --- while many of today's problem characters suffer too little inhibition, not too much fear: a character disturbance, not a neurosis.
  \item The manipulator generally knows what he is doing; what he lacks is restraint, not insight --- so sympathetic interpretation and insight-work miss him entirely.
  \item The victim profile: feeling confused and crazy, distrusting without pinpointing why, getting angry and ending up guilty, confronting and winding up on the defensive, saying yes when meaning no.
  \item Everyday covert-aggressors are the common case --- workmates, associates, relatives --- not the headline wolves.
  \item The manipulative child chapter extends the category developmentally: children learn (or are never taught) to discipline aggressive instincts; permissiveness, indulgence and lack of accountability produce more covert-aggressive children.
\end{distilled}
\begin{unique}
  \item The covert-aggressive category itself, the character-disturbance-versus-neurosis reframing, and the developmental account in the manipulative-child chapter are unique to B6 in this corpus.
\end{unique}
\begin{terms}
  \item \emph{covert aggression} --- B6's term for a disguised fight: pursuing one's wants and resisting others' while concealing the aggressive intent.
  \item \emph{character disturbance} --- B6's framing for deficient restraint over instinct, as against neurotic over-inhibition.
\end{terms}
